\section{Experiments and Results}

\paragraph{Experimental protocol.}

All effectiveness experiments and system comparisons reported in this
section use the 22-topic development set; the official 119-topic data
are used only for structural submission diagnostics. Retrieval uses
three projected UMBRELA qrels, relevance threshold
$\mathrm{rel}\geq2$, and linear-gain nDCG. P@10, Recall, MAP, and MRR
use binarized $\mathrm{rel}\geq2$ judgments. For each Retrieval run, we
evaluate independently against the three projected UMBRELA qrels and
report the arithmetic mean of the resulting aggregate metrics. The
qrels are not merged. We validate rank continuity, duplicate
identifiers, JSON structure, citation indices, and response-length
compliance before submission.

\paragraph{Official submission.}

Our team (team ID
\texttt{2026 cfda rag}) submitted official runs for the 119-topic test
set. Team submission records list \texttt{cfda-vfs-deep} (priority 1),
corresponding to the selected deep-tail retrieval configuration, and
\texttt{cfda-vfs-unc} (priority 2), corresponding
to the selected pre-deep-tail retrieval configuration described in
\S\ref{sec:retrieval-method}, as the two Retrieval runs. They list
\texttt{cfda-w5c} (priority 1) as the RAG run, corresponding to the
selected bounded 60-document acquisition configuration described in
\S\ref{sec:rag-method}.
Both Retrieval submissions were serialized with the
variable-depth rule of Eq.~\ref{eq:vark} at $K=5{,}000$, giving
per-topic depths from 137 to 4,218 (median 1,080); each run file
contains 176,224 rows.
The two runs share the same underlying candidate pool, unchanged
top-100 prefix, and per-topic depth. Deep-tail reranking changes the tail
order for every topic and, for 103 of 119 topics, changes which
documents fall inside the submitted prefix. At $K=1{,}000$, an
intermediate official serialization made 63 of 119 topics hit the cap.
Because the task specifies no maximum depth and the pool held up to
5,000 documents, we retained the development-selected $\tau=0.20$ and
rank-1 anchor and raised only the team cap to $K=5{,}000$, reducing cap
domination. This decision used structural diagnostics, not official
judgments. Official
test-set scores are assigned by the track
organizers after submission and are not yet available at the time of
this draft; all scores reported below are development-set diagnostics,
as noted throughout.

The selected Retrieval output is the variable-depth system. The fixed
top-1000 output is retained as a fixed-depth comparison because it uses a
different submission depth: its Recall@1000 of 0.3159 is not substituted
for the variable-depth Recall@submitted-$k$ of 0.2343.

For RAG, we compare the 12-document-cap baseline, the 60-document
deep-reading configuration, and the selected bounded 60-document
acquisition configuration; extended reading is reserved for failure
analysis. For $V_{\mathrm{strict}}$, the selected-minus-deep-reading
mean paired topic-level difference is $-0.0042$, with an approximate
95\% paired $t$ interval of $[-0.0497,+0.0412]$ and an 8/11/3
win/loss/tie record. Selection reflects stronger citation-support
behavior; the $V_{\mathrm{strict}}$ difference is inconclusive, and the stored
variants do not support a single-factor causal interpretation.

Key implementation and operating parameters for both systems are summarized in Appendix
Table~\ref{tab:appendix-retrieval-parameters} (Retrieval) and
Table~\ref{tab:appendix-rag-parameters} (RAG).
